problem:  
Let \( G \) be a compact simple Lie group and \( P \subset G \) a parabolic subgroup, with \( M = G/P \) the associated generalized flag manifold. Denote by  
\[
\mathfrak{g} = \mathfrak{p} \oplus \mathfrak{m} = \mathfrak{p} \oplus \bigoplus_{i=1}^r \mathfrak{m}_i
\]
the isotropy decomposition into irreducible \( \operatorname{Ad}(P) \)-submodules.  

A \( G \)-invariant Riemannian metric on \( M \) is determined (up to scale) by positive parameters \( (x_1, \dots, x_r) \), where \( x_i \) rescales the negative Killing form on each \(\mathfrak{m}_i\). The Einstein condition for invariance yields a system of algebraic equations
\[
F_i(x_1,\dots,x_r) = \lambda x_i,\quad i=1,\dots,r,
\]
whose solutions correspond to invariant Einstein metrics on \( M \). The canonical Kähler–Einstein metric \(g_{KE}\) corresponds to a specific point \(x^{KE} \in (\mathbb{R}_{>0})^r\).

**Problem:**  
Characterize, in terms of the isotropy representation of \(G/P\), the Lie algebra decomposition, and root system data, the conditions under which the Jacobian matrix \(J(x^{KE}) = \left( \frac{\partial F_i}{\partial x_j}(x^{KE}) - \lambda \delta_{ij} \right)\) of the invariant Einstein system becomes singular at the Kähler–Einstein point.  

In particular, does there exist any flag manifold \( M = G/P \) for which \( \det J(x^{KE}) = 0 \)?  
If such degeneracy occurs, what representation-theoretic mechanism forces the Einstein operator to have zero modes—e.g., multiplicities or symmetry coincidences among the isotropy summands arising from the root system structure of \(G\)?

Develop Lie-theoretic or combinatorial criteria predicting when the canonical Kähler–Einstein metric is *non-isolated* in the space of invariant Einstein metrics.

---

Why is it a "good" problem:  
This problem targets the **structural core** of homogeneous Einstein geometry. It asks whether the canonical Kähler–Einstein metric on a flag manifold can ever be a *degenerate solution* of the invariant Einstein system, and if so, how that degeneracy reflects the underlying Lie algebraic geometry (root combinatorics, isotropy representation multiplicities, or painting symmetries in the Dynkin diagram).

The question is conceptually simple—does the Jacobian of the Einstein system ever vanish?—yet its implications are deep. A positive result would reveal hidden moduli within the apparently rigid lattice of homogeneous Einstein metrics, while a negative one would confirm a profound structural rigidity principle for all flag manifolds. Either outcome would sharpen our understanding of the interplay between the algebraic data of \(G/P\) and the analytic geometry of the Einstein equation.

This problem sits squarely at the intersection of **algebraic representation theory**, **differential geometry**, and **Einstein analysis**, without contradicting known facts about invariant complex structures. Its resolution would clarify whether the algebraic Einstein equations can encode genuine geometric degeneracies—a fundamentally open and structurally significant question.